% ============================================================================
%  บทที่ 5: ผลการทดสอบและการทวนสอบ
% ============================================================================

\chapter{ผลการทดสอบและการทวนสอบ}
\label{ch:results}

% ============================================================================
%  บทที่ 5: ผลการทดสอบและการทวนสอบ
%  Section 5.1 — ผลการทดสอบระบบย่อย (ปรับปรุง)
% ============================================================================

% ============================================================================
\section{ผลการทดสอบระบบย่อย}
\label{sec:subsystem_test}
% ============================================================================

การทดสอบระบบย่อยดำเนินการก่อนการทดสอบสมรรถนะเต็มรูปแบบ
เพื่อยืนยันว่าส่วนประกอบแต่ละชิ้นทำงานถูกต้องตามข้อกำหนดทางกลและทางไฟฟ้า
ก่อนนำไปประกอบเป็นระบบควบคุมแบบ Closed-loop
ข้อมูลในหัวข้อนี้เก็บผ่าน STM32CubeMonitor โดย Log ตัวแปรสถานะผ่าน
SWD/JTAG ที่ Sample Rate 100\,Hz เว้นแต่ระบุเป็นอื่น

% ----------------------------------------------------------------------------
\subsection{การวัดระยะคลอน (Backlash Measurement)}
\label{subsec:backlash}
% ----------------------------------------------------------------------------

\subsubsection{วัตถุประสงค์และความสำคัญ}

ระยะคลอน (Backlash) ในระบบส่งกำลังทำให้เกิด Dead-zone ในลูปควบคุม
ซึ่งส่งผลโดยตรงต่อความแม่นยำตำแหน่ง (A1: $\leq \pm 0.5\degree$)
และความซ้ำซากได้ (A2: $\leq \pm 0.1\degree$)
ในระบบนี้มีแหล่งกำเนิดระยะคลอนสองแหล่งหลัก ได้แก่
Planetary Gearbox (อัตราทด $N_g = 24$) และ Belt Drive (อัตราทด $N_b = 70/24$)
การวัดจึงแยกเป็นสองจุดเพื่อระบุแหล่งที่มาของความคลาดเคลื่อน

\subsubsection{วิธีการทดสอบ}

ล็อกตำแหน่งแขนกลที่มุม $\theta = 0\degree$, $40\degree$, $90\degree$
ด้วยการให้ Controller ส่ง PWM = 0 และ Brake Enable
จากนั้นออกแรงกดที่ปลายแขน (ระยะ $r = [VAL\_ARM\_LENGTH]$\,mm จากศูนย์กลาง)
ทั้งในทิศทาง $+$ และ $-$ และอ่านค่า Encoder Count
ระยะคลอนคำนวณจาก:

\begin{equation}
    \delta_{backlash} = \frac{\Delta \text{counts}}{8192 \cdot N_{total}} \times 360\degree
    \label{eq:backlash_calc}
\end{equation}

โดย $N_{total} = 70$ คืออัตราทดรวมของระบบส่งกำลัง
และ 8192 counts/rev คือ Resolution ของ Encoder AMT-103
ที่ $\times 4$ QEI บน STM32G474RE

% -----------------------------------------------------------------------
% รูปที่: Backlash Measurement Setup
% ชื่อไฟล์แนะนำ: fig_backlash_setup.png
% เนื้อหา: ภาพถ่ายหรือ Diagram แสดงจุดที่ออกแรงกดที่ปลายแขน
%          ระบุทิศทาง +/- และตำแหน่ง Encoder readout
% -----------------------------------------------------------------------
% \begin{figure}[H]
%     \centering
%     \includegraphics[width=0.65\textwidth]{fig_backlash_setup}
%     \caption{Setup การวัดระยะคลอน: ออกแรงกดที่ปลายแขนทั้งสองทิศทาง
%              และบันทึกค่า Encoder ผ่าน STM32CubeMonitor}
%     \label{fig:backlash_setup}
% \end{figure}

\subsubsection{ผลการวัด}

\begin{table}[H]
  \centering
  \caption{ผลการวัดระยะคลอนของระบบส่งกำลัง
           (ค่าเฉลี่ยจากการวัด 3 ตำแหน่ง $\times$ 5 ครั้ง)}
  \label{tab:backlash}
  \small
  \begin{tabular}{>{\raggedright\arraybackslash}p{3.5cm} >{\raggedright\arraybackslash}p{3.5cm} >{\raggedright\arraybackslash}p{3.5cm} >{\raggedright\arraybackslash}p{3.0cm}}
    \toprule
    จุดวัด & ระยะคลอน ($\degree$) & $\sigma$ ($\degree$) & หมายเหตุ \\
    \midrule
    Planetary Gearbox & 0.048 & 0.003 & วัดที่ Output shaft กล่อง \\
    Belt Drive        & 0.095 & 0.008 & วัดที่ Output shaft แขน  \\
    \midrule
    รวมระบบ           & 0.143 & -- &
      เกณฑ์ A2: $\leq \pm 0.1\degree$ \quad ไม่ผ่าน \\
    \bottomrule
  \end{tabular}
\end{table}

% -----------------------------------------------------------------------
% รูปที่: Backlash Encoder Trace
% ชื่อไฟล์แนะนำ: fig_backlash_trace.png
% เนื้อหา: Plot theta_arm vs time จาก STM32CubeMonitor
%          แสดงช่วง Dead-zone ขณะเปลี่ยนทิศทางแรง
%          Annotate: delta_counts ที่ใช้คำนวณ backlash
% -----------------------------------------------------------------------
% \begin{figure}[H]
%     \centering
%     \includegraphics[width=0.75\textwidth]{fig_backlash_trace}
%     \caption{Encoder Trace ระหว่างการวัดระยะคลอน
%              แสดงช่วง Dead-zone เมื่อเปลี่ยนทิศทางการออกแรง}
%     \label{fig:backlash_trace}
% \end{figure}

% ----------------------------------------------------------------------------
\subsection{การทดสอบ Encoder และ Sensor}
\label{subsec:sensor_test}
% ----------------------------------------------------------------------------

\subsubsection{การทดสอบ Encoder AMT-103}
\label{subsubsec:encoder_test}

\paragraph{วัตถุประสงค์}
Encoder AMT-103 เป็นอุปกรณ์หลักในการ Feedback ตำแหน่งและความเร็ว
การยืนยัน Resolution, Noise Floor, และการทำงานของ Index Pulse
มีผลโดยตรงต่อ Bandwidth ของ Velocity Loop และความแม่นยำของ Position Controller

\paragraph{การทดสอบ Resolution}
หมุนแขนกลผ่านมุม $360\degree$ หนึ่งรอบเต็ม (ที่ Output Shaft)
แล้วนับ Encoder Counts ที่ STM32 อ่านได้ผ่าน QEI (Quadrature Encoder Interface)
ค่าที่คาดหวังคือ $8192 \times N_{total} = 8192 \times 70 = 573{,}440$ counts/rev
ที่ Output Shaft (แขน)
หรือเทียบได้กับ $\approx 22.76$ counts/$\degree$ ที่ Output Shaft

\paragraph{การทดสอบ Noise Floor}
ล็อกแขนกลที่ตำแหน่ง $\theta = 0\degree$ (Brake Enable, PWM = 0)
เก็บข้อมูล Encoder Count ต่อเนื่อง 500 Samples ที่ 1\,kHz
ผ่าน STM32CubeMonitor บันทึก \texttt{encoder\_raw}
และคำนวณ Peak-to-peak Noise เพื่อประเมินว่า Noise
ส่งผลต่อ Velocity Estimate ของ Kalman Filter อย่างไร

% -----------------------------------------------------------------------
% รูปที่: Encoder Noise Floor
% ชื่อไฟล์แนะนำ: fig_encoder_noise.png
% เนื้อหา: 2 subplots:
%   (บน) encoder_raw counts vs time (500 samples, standstill)
%   (ล่าง) Histogram ของ deviation จากค่าเฉลี่ย
%          Annotate: peak-to-peak range
% -----------------------------------------------------------------------
\begin{figure}[H]
    \centering
    \includegraphics[width=0.80\textwidth]{fig_encoder_noise}
    \caption{Encoder Noise Floor ขณะหยุดนิ่ง:
             Time Series (บน) และ Histogram ของ Deviation (ล่าง)}
    \label{fig:encoder_noise}
\end{figure}

\begin{table}[H]
  \centering
  \caption{ผลการทดสอบ Encoder AMT-103}
  \label{tab:encoder_test}
  \small
  \begin{tabular}{>{\raggedright\arraybackslash}p{6.0cm} >{\centering\arraybackslash}p{2.5cm} >{\centering\arraybackslash}p{3.0cm} >{\centering\arraybackslash}p{2.0cm}}
    \toprule
    รายการทดสอบ & ผล & เกณฑ์ & สถานะ \\
    \midrule
    Resolution (counts/rev ที่ Motor Shaft) & 8192 & 8192 ($\times 4$ QEI) & ผ่าน \\
    ค่าคงที่ต่อ $1\degree$ ที่ Output Shaft         & 1592.9 & $\approx 22.76$ & ผ่าน \\
    Noise Peak-to-peak (ขณะหยุดนิ่ง)       & $<1$ & $\leq 2$ counts & ผ่าน \\
    Index Pulse ทำงานถูกต้อง                & ผ่าน & ผ่าน & ผ่าน \\
    \bottomrule
  \end{tabular}
\end{table}

\subsubsection{การทดสอบ Proximity Sensor (Homing)}
\label{subsubsec:prox_test}

\paragraph{วัตถุประสงค์}
Proximity Sensor ใช้กำหนดจุด Home ($\theta = 0\degree$) ทุกครั้งที่เริ่มต้นระบบ
ความซ้ำซากได้ของ Homing Position มีผลโดยตรงต่อ Absolute Accuracy (A1)
เนื่องจาก Encoder เป็น Incremental Type
ซึ่งหมายความว่าข้อผิดพลาดของ Homing จะสะสมเป็น Offset ถาวร
ตลอดทั้ง Cycle การทำงาน

\paragraph{วิธีการทดสอบ}
ดำเนิน Homing Sequence 10 ครั้งติดต่อกัน โดยในแต่ละครั้ง:
\begin{enumerate}
    \item หมุนแขนออกจาก Home ไปยังตำแหน่งสุ่ม $\theta \in [30\degree, 90\degree]$
    \item สั่ง Homing Sequence: แขนหมุนกลับด้วยความเร็ว Slow ($\omega = 0.5$\,rad/s)
          จนกว่า Proximity Sensor จะ Trigger
    \item บันทึกค่า Encoder ณ จุดที่ Sensor Trigger ผ่าน STM32CubeMonitor (\texttt{theta\_home\_raw})
\end{enumerate}
ความซ้ำซากได้คำนวณจาก $\sigma$ และ Max deviation ของ 10 ค่า

% -----------------------------------------------------------------------
% รูปที่: Homing Repeatability
% ชื่อไฟล์แนะนำ: fig_homing_repeat.png
% เนื้อหา: Scatter plot แสดง theta_home ทั้ง 10 ครั้ง
%          เส้น dashed = ค่าเฉลี่ย
%          Shaded band = ±0.1\degree (เกณฑ์ A2)
% -----------------------------------------------------------------------
\begin{figure}[H]
    \centering
    \includegraphics[width=0.70\textwidth]{fig_homing_repeat}
    \caption{ความซ้ำซากได้ของ Homing Position 10 ครั้ง
             เทียบกับเกณฑ์ $\pm 0.1\degree$ (A2)}
    \label{fig:homing_repeat}
\end{figure}

\begin{table}[H]
  \centering
  \caption{ผลการทดสอบ Proximity Sensor สำหรับ Homing}
  \label{tab:prox_test}
  \small
  \begin{tabular}{>{\raggedright\arraybackslash}p{4.5cm} >{\centering\arraybackslash}p{3.0cm} >{\centering\arraybackslash}p{3.0cm} >{\centering\arraybackslash}p{3.0cm}}
    \toprule
    รายการทดสอบ & ผล & เกณฑ์ & สถานะ \\
    \midrule
    Switching Distance                    & -- & -- & -- \\
    ความซ้ำซากได้ของ Homing ($\sigma$)    & 0.0000$\degree$ & -- & -- \\
    Max Deviation (10 reps)               & 0.0000$\degree$ & $\leq \pm 0.1\degree$ & ผ่าน \\
    \bottomrule
  \end{tabular}
\end{table}

% ----------------------------------------------------------------------------
\subsection{การทดสอบ Motor Driver (Cytron MD10C)}
\label{subsec:motor_driver_test}
% ----------------------------------------------------------------------------

\subsubsection{วัตถุประสงค์}

Cytron MD10C รับสัญญาณ PWM จาก STM32G474RE และขับ RS-775
ที่แรงดัน 24\,V/10\,A
การทดสอบ Current Step Response ยืนยันว่า Driver สามารถตามทัน
Current Command จาก Velocity Loop ซึ่งทำงานที่ 1\,kHz
โดยไม่เกิด Oscillation หรือ Delay ที่ส่งผลต่อ Inner Loop Bandwidth

\subsubsection{วิธีการทดสอบ}

ให้ Position Loop ส่ง Step Command $\theta_{ref} = 40\degree$ จากสถานะนิ่ง
ขณะที่ Log ตัวแปรต่อไปนี้ผ่าน STM32CubeMonitor ที่ Sample Rate 1\,kHz:

\begin{itemize}
    \item \texttt{pwm\_cmd} — PWM Duty Cycle (\%) ที่ส่งไปยัง MD10C
    \item \texttt{ia\_measured} — กระแส Armature (A) ที่ประมาณจาก Kalman Filter
    \item \texttt{omega\_arm} — ความเร็วเชิงมุม (rad/s) ที่ Output Shaft
\end{itemize}

\subsubsection{ผลการทดสอบ}

\noindent ข้อมูล Motor Driver Response ไม่สามารถเก็บได้จากการทดสอบนี้
เนื่องจาก PWM column ไม่ถูก Export ผ่าน Dashboard Log
การยืนยันทำผ่านการสังเกตการทำงานจริงของ Motor ขณะ Cycle Test

% -----------------------------------------------------------------------
% รูปที่: Motor Driver Step Response
% ชื่อไฟล์แนะนำ: fig_driver_step_response.png
% เนื้อหา: 3 subplots (จาก STM32CubeMonitor log):
%   (บน)   pwm_cmd (%) vs time
%   (กลาง) ia_measured (A) vs time — annotate i_peak
%   (ล่าง) omega_arm (rad/s) vs time
%          x-axis: 0–2s, ครอบคลุม Transient ช่วงแรก
% -----------------------------------------------------------------------
% \begin{figure}[H]
%     \centering
%     \includegraphics[width=0.88\textwidth]{fig_driver_step_response}
%     \caption{Motor Driver Response ต่อ Step Position Command $40\degree$:
%              PWM Command, กระแส Armature, และความเร็วเชิงมุม}
%     \label{fig:driver_step_response}
% \end{figure}

% ----------------------------------------------------------------------------
\subsection{การทดสอบ E-Stop Response Time}
\label{subsec:estop_test}
% ----------------------------------------------------------------------------

\subsubsection{วัตถุประสงค์}

E-Stop (Emergency Stop) ตัดกำลังไฟให้ Cytron MD10C ผ่าน Contactor
โดย 5\,V Logic Bus และ STM32 ยังคงมีไฟเลี้ยงอยู่ (ข้อกำหนด E6)
การทดสอบนี้วัดเวลาตั้งแต่กดปุ่ม E-Stop จนกระทั่ง $\omega_{arm} \approx 0$
เพื่อประเมินความปลอดภัยของระบบภายใต้สถานการณ์ฉุกเฉิน

\subsubsection{วิธีการทดสอบ}

\begin{enumerate}
    \item สั่งให้แขนเคลื่อนที่ด้วยความเร็วค่าหนึ่ง
    \item กดปุ่ม E-Stop ขณะแขนกำลังเคลื่อนที่
    \item Log ตัวแปร \texttt{omega\_arm}, \texttt{pwm\_cmd}, \texttt{estop\_flag}
          ผ่าน STM32CubeMonitor ที่ 1\,kHz
    \item คำนวณ $t_{stop}$ = เวลาตั้งแต่ \texttt{estop\_flag} = 1
          จนถึง $|\omega_{arm}| < 0.01$\,rad/s
    \item ทำซ้ำ 5 ครั้ง
\end{enumerate}

\subsubsection{ผลการทดสอบ}

\noindent การทดสอบ E-Stop ยืนยันว่าการกดปุ่มฉุกเฉินตัดกำลัง Motor ทันที
และ STM32 Controller ยังคงทำงานต่อเนื่อง (ข้อกำหนด E6)
ตัวเลข $t_{stop}$ เชิงปริมาณอยู่ระหว่างรอการเก็บข้อมูลเพิ่มเติม

% -----------------------------------------------------------------------
% รูปที่: E-Stop Response Time Trace
% ชื่อไฟล์แนะนำ: fig_estop_trace.png
% เนื้อหา: 3 subplots (จาก STM32CubeMonitor):
%   (บน)   estop_flag (0/1) vs time — แสดงจุดที่กด E-Stop
%   (กลาง) pwm_cmd (%) vs time — ตัดเป็น 0 ทันที
%   (ล่าง) omega_arm (rad/s) vs time
%          Annotate: t_stop = เวลาหยุดสนิท
%          Annotate: Controller ยังทำงานอยู่ (log ต่อเนื่อง)
% -----------------------------------------------------------------------
% \begin{figure}[H]
%     \centering
%     \includegraphics[width=0.80\textwidth]{fig_estop_trace}
%     \caption{E-Stop Response: E-Stop Flag, PWM Command, และ
%              ความเร็วเชิงมุม $\omega_{arm}$ แสดงการหยุดของแขนกล
%              หลังกด E-Stop (ข้อกำหนด E6)}
%     \label{fig:estop_trace}
% \end{figure}

% ----------------------------------------------------------------------------
\subsection{การทดสอบการสื่อสาร Modbus RTU}
\label{subsec:modbus_test}
% ----------------------------------------------------------------------------

\subsubsection{วัตถุประสงค์}

ระบบรับ Command จาก Base System ผ่าน Modbus RTU (LPUART1, 230400\,baud, 8N1)
โดย STM32G474RE ทำหน้าที่เป็น Slave (Address 21)
การทดสอบยืนยันว่า Function Code FC03 (Read Holding Registers)
และ FC06 (Write Single Register) ทำงานถูกต้อง
และ Round-trip Latency เหมาะสมกับ Cycle Time ที่กำหนด ($\leq 35$\,s)

\subsubsection{วิธีการทดสอบ}

ใช้ Python Script ส่ง Modbus Request จาก Base System PC ไปยัง STM32
แล้ววัด Round-trip Time โดย Timestamp ทั้งฝั่ง Master (Python) และ Slave (STM32)
\begin{itemize}
    \item \textbf{FC03}: อ่านค่า \texttt{theta\_arm}, \texttt{omega\_arm}, \texttt{system\_status}
          จำนวน 10 Registers
    \item \textbf{FC06}: เขียน \texttt{position\_cmd} (1 Register) พร้อมตรวจสอบว่า
          STM32 ตอบรับและ Execute ถูกต้อง
    \item ทดสอบ 100 Transactions ต่อเนื่อง วัด Latency และ Error Rate
\end{itemize}

\subsubsection{ผลการทดสอบ}

\noindent Modbus RTU FC03 และ FC06 ทำงานถูกต้องโดยไม่พบ CRC Error
ตลอดการทดสอบ Full Cycle Run จำนวน 3 รอบ
ค่า Round-trip Latency เชิงปริมาณอยู่ระหว่างรอการเก็บข้อมูล

% -----------------------------------------------------------------------
% รูปที่: Modbus Latency Distribution
% ชื่อไฟล์แนะนำ: fig_modbus_latency.png
% เนื้อหา: Histogram ของ Round-trip latency (100 transactions)
%          แยก FC03 (น้ำเงิน) และ FC06 (แดง)
%          เส้น dashed = ค่าเฉลี่ย, เส้น dotted = Max
% -----------------------------------------------------------------------
% \begin{figure}[H]
%     \centering
%     \includegraphics[width=0.72\textwidth]{fig_modbus_latency}
%     \caption{การกระจาย Round-trip Latency ของ Modbus RTU
%              FC03 และ FC06 จาก 100 Transactions}
%     \label{fig:modbus_latency}
% \end{figure}
% ============================================================================
\section{ผลการระบุระบบและการจำลอง}
\label{sec:sysid_results}
% ============================================================================

หัวข้อนี้รายงานผลการระบุพารามิเตอร์ของมอเตอร์และแบบจำลองระบบ
ซึ่งเป็นรากฐานสำคัญสำหรับการออกแบบ Cascade PID และ Kalman Filter
กระบวนการระบุระบบแบ่งออกเป็นสองขั้นตอน ได้แก่
การวัดพารามิเตอร์วงจรไฟฟ้าจาก Locked Rotor Test
และการระบุพารามิเตอร์เชิงกลจาก Chirp Excitation Test
รายละเอียดทฤษฎีและขั้นตอนการออกแบบแสดงไว้ใน \cref{ch:design}

% ----------------------------------------------------------------------------
\subsection{ผลการระบุพารามิเตอร์ของมอเตอร์}
\label{subsec:sysid_params_result}
% ----------------------------------------------------------------------------

\subsubsection{ผล Locked Rotor Test}
\label{subsubsec:locked_rotor_result}

การวัด $R_a$ และ $L_a$ ดำเนินการก่อนการระบุพารามิเตอร์เชิงกล
เนื่องจากค่าทั้งสองเป็น Input ที่ทราบค่าแน่นอนสำหรับแบบจำลองทางไฟฟ้า
หากนำ $R_a$ และ $L_a$ ไปประมาณร่วมกับ $K_t$, $K_b$, $B$, $J$
จะทำให้ Optimization Problem มีระดับความอิสระสูงเกินไปและเกิด Parameter Correlation
การล็อกโรเตอร์ตัด Back-EMF ออกจากวงจร ($\omega = 0 \Rightarrow e = K_b\omega = 0$)
ทำให้สามารถระบุ $R_a = V_{ss}/I_{ss}$ และ $L_a = R_a \tau_{RL}$
ได้อย่างเป็นอิสระจากพารามิเตอร์เชิงกล

\begin{table}[H]
  \centering
  \caption{ผลการวัด $R_a$ และ $L_a$ จาก Locked Rotor Test
           ที่ $V_{in} = 12$\,V (3 runs)}
  \label{tab:locked_rotor}
  \small
  \begin{tabular}{>{\centering\arraybackslash}p{3.3cm} >{\centering\arraybackslash}p{3.4cm} >{\centering\arraybackslash}p{3.4cm} >{\centering\arraybackslash}p{3.4cm}}
    \toprule
    Run & $R_a$ ($\Omega$) & $L_a$ (mH) & หมายเหตุ \\
    \midrule
    1 & 1.5216 & 1.4595 & -- \\
    2 & 1.4385 & 1.3846 & -- \\
    3 & 1.4000 & 1.5000 & -- \\
    \midrule
    เฉลี่ย & \textbf{1.4534} & \textbf{1.448} & ค่าที่ใช้ใน Design \\
    $\sigma$ & 0.0523 & 0.000055 & -- \\
    \bottomrule
  \end{tabular}
\end{table}

ค่า $\sigma$ ที่ต่ำบ่งชี้ว่าการวัดมีความสม่ำเสมอดี
ค่าคงเวลาทางไฟฟ้า $\tau_{RL} = L_a/R_a = 1.448/1453.4 \approx 1.0$\,ms
สอดคล้องกับ Velocity Loop Sampling Period ที่ 1\,kHz
ยืนยันว่า Current Dynamics เร็วพอที่จะถูกมองเป็น Static Gain
จากมุมมองของ Velocity Loop ได้

% -----------------------------------------------------------------------
% รูปที่: Locked Rotor Oscilloscope Waveform
% ชื่อไฟล์แนะนำ: fig_locked_rotor_scope.png
% เนื้อหา: Oscilloscope waveform แสดง V_in step และ i_a response
%          Annotate: V_ss, I_ss สำหรับคำนวณ R_a
%          Annotate: tau_RL = L_a/R_a (จุด 63.2% ของ I_ss)
% -----------------------------------------------------------------------
% \begin{figure}[H]
%     \centering
%     \includegraphics[width=0.70\textwidth]{fig_locked_rotor_scope}
%     \caption{Oscilloscope Waveform จาก Locked Rotor Test ($V_{in} = 12$\,V)
%              แสดง Steady-state Current $I_{ss}$ และค่าคงเวลา $\tau_{RL}$
%              ที่ใช้คำนวณ $R_a$ และ $L_a$}
%     \label{fig:locked_rotor_scope}
% \end{figure}

\subsubsection{ผล Parameter Estimation จาก Chirp Test}
\label{subsubsec:chirp_result}

การระบุพารามิเตอร์เชิงกล ($K_t$, $K_b$, $B$, $J$, $\eta$) ใช้ Chirp Excitation
ที่ความถี่ 0.1--2.0\,Hz เพื่อครอบคลุมย่านความถี่ที่น่าสนใจของ Velocity Loop (DC ถึง ~5\,Hz)
โดยไม่กระตุ้น Resonance ของ Connection Rod ($f_{n,rod} = 12.13$\,Hz)

\paragraph{การ Preprocessing สัญญาณ}
ก่อนนำข้อมูลเข้า Least-Squares Estimator จำเป็นต้องกรองสัญญาณ
เพื่อลด Encoder Quantization Noise และ PWM Switching Noise (19.6\,kHz)
โดยใช้ Butterworth Filter อันดับ 4 ที่ความถี่ตัด $f_c = 10$\,Hz
ด้วยวิธี Zero-phase Filtering (\texttt{filtfilt}) เพื่อไม่ให้เกิด Phase Shift
ที่จะทำให้ Input--Output Causality ผิดเพี้ยน
การเลือก $f_c = 10$\,Hz มีเหตุผลดังนี้:
\begin{itemize}
    \item \textbf{ต่ำกว่า $f_{n,rod}$}: หลีกเลี่ยงการปนเปื้อนของ Rod Dynamics
          เข้ามาใน Motor Model
    \item \textbf{สูงกว่า Band ของ Chirp}: ช่วง 0.1--2\,Hz อยู่ต่ำกว่า $f_c$
          มาก ทำให้สัญญาณที่ต้องการผ่านโดยไม่ถูก Attenuate
    \item \textbf{ยืนยันด้วย ETFE}: Bode Plot จาก Empirical Transfer Function Estimate
          (ETFE) ก่อนและหลัง Filter ให้ผลสอดคล้องกันในย่าน 0.1--2\,Hz
          ยืนยันว่า Filter ไม่บิดเบือนพารามิเตอร์ที่ระบุ
\end{itemize}

% -----------------------------------------------------------------------
% รูปที่: Chirp Raw vs Filtered
% ชื่อไฟล์แนะนำ: fig_chirp_raw_filtered.png
% เนื้อหา: 2 subplots:
%   (บน) V_in raw vs filtered vs time
%   (ล่าง) omega_arm raw vs filtered vs time
%   แสดงช่วง Chirp ความถี่ต่ำ (0.1 Hz) และสูง (2 Hz)
%   Annotate: f_c = 10 Hz cutoff ใน legend
% -----------------------------------------------------------------------
% \begin{figure}[H]
%     \centering
%     \includegraphics[width=0.88\textwidth]{fig_chirp_raw_filtered}
%     \caption{สัญญาณ Chirp ก่อนและหลังกรองด้วย Butterworth อันดับ~4
%              ที่ $f_c = 10$\,Hz ด้วยวิธี Zero-phase (\texttt{filtfilt})
%              แสดง $V_{in}$ (บน) และ $\omega_{arm}$ (ล่าง)}
%     \label{fig:chirp_raw_filtered}
% \end{figure}

\paragraph{ผลการประมาณพารามิเตอร์}
ดำเนิน Chirp Test จำนวน 9 Experiments ที่แอมพลิจูดและช่วงความถี่ต่างกัน
เพื่อประเมิน Uncertainty ของแต่ละพารามิเตอร์
ค่าเฉลี่ยและ Standard Deviation จากทั้ง 9 Experiments แสดงใน\cref{tab:chirp_results}

\begin{table}[H]
  \centering
  \caption{ผลการ Estimate พารามิเตอร์จาก Chirp Test (สรุปจาก 9 Experiments)}
  \label{tab:chirp_results}
  \small
  \begin{tabular}{>{\raggedright\arraybackslash}p{2.6cm} >{\centering\arraybackslash}p{1.2cm} >{\centering\arraybackslash}p{1.2cm} >{\centering\arraybackslash}p{2.0cm} >{\centering\arraybackslash}p{2.0cm} >{\centering\arraybackslash}p{1.8cm} >{\centering\arraybackslash}p{1.7cm}}
    \toprule
    พารามิเตอร์ & สัญลักษณ์ & เฉลี่ย & $\sigma$ & Uncertainty (\%) & หน่วย & หมายเหตุ \\
    \midrule
    Back-EMF Constant & $K_b$  & 0.04165 & 0.00201  & 4.86  & V$\cdot$s/rad  & -- \\
    Torque Constant   & $K_t$  & 0.04065 & 0.00356  & 8.41  & N$\cdot$m/A    & -- \\
    Viscous Damping   & $B$    & 0.19279 & 0.01658   & 8.71   & N$\cdot$m$\cdot$s/rad & สูงสุด \\
    Moment of Inertia & $J$    & 0.72762 & 0.35565   & 38.64   & kg$\cdot$m$^2$ & -- \\
    Efficiency        & $\eta$ & 0.836   & 0.01295 & 1.56 & --             & -- \\
    \bottomrule
  \end{tabular}
\end{table}

พารามิเตอร์ $B$ (Viscous Damping) มักมี Uncertainty สูงที่สุดในกลุ่ม
เนื่องจาก Damping จริงประกอบด้วยทั้ง Viscous, Coulomb Friction และ Stiction
ซึ่งแบบจำลอง Linear Viscous ไม่สามารถจับได้ครบ
อย่างไรก็ตาม Kalman Filter ออกแบบให้ประมาณ Disturbance Torque $\tau_d$
แบบ Real-time จึงชดเชย Unmodeled Friction ได้โดยตรง
ทำให้ Uncertainty ของ $B$ มีผลต่อ Performance น้อยกว่าที่คาดในทางทฤษฎี

% -----------------------------------------------------------------------
% รูปที่: ETFE Bode Plot Comparison
% ชื่อไฟล์แนะนำ: fig_etfe_bode.png
% เนื้อหา: Bode Plot (Magnitude + Phase) 2 เส้น:
%   เส้นสีเทา = ETFE จากข้อมูล Raw
%   เส้นสีน้ำเงิน = ETFE จากข้อมูล Filtered
%   เส้นสีแดง dashed = Identified Model G(jω)
%   Annotate: f_c = 10 Hz, f_n_rod = 12.13 Hz
%   แสดงว่า Filter ไม่ distort ย่าน 0.1–2 Hz
% -----------------------------------------------------------------------
% \begin{figure}[H]
%     \centering
%     \includegraphics[width=0.88\textwidth]{fig_etfe_bode}
%     \caption{ETFE Bode Plot เปรียบเทียบ Raw, Filtered และ Identified Model
%              ยืนยันว่า Butterworth Filter ไม่บิดเบือนพารามิเตอร์
%              ในย่านความถี่ระบุ (0.1--2\,Hz)}
%     \label{fig:etfe_bode}
% \end{figure}

% ----------------------------------------------------------------------------
\subsection{ผลการตรวจสอบความถูกต้องของแบบจำลอง (Model Validation)}
\label{subsec:model_validation_results}
% ----------------------------------------------------------------------------

การตรวจสอบความถูกต้องของแบบจำลองดำเนินการด้วยวิธี Cross-Validation
โดยใช้ชุดข้อมูล Chirp ที่ความถี่ต่างกัน 3 ชุด (0.5, 1.0, 2.0\,Hz)
ซึ่งเป็นชุดข้อมูลที่ \emph{ไม่ได้ใช้} ใน Estimation
เพื่อยืนยันว่าแบบจำลองสามารถ Generalize ได้จริง
ไม่ใช่เพียง Overfit กับข้อมูลที่ใช้ระบุ

Fit Score คำนวณจาก:
\begin{equation}
    \text{Fit} = \left(1 - \frac{\|\hat{y} - y\|}{\|y - \bar{y}\|}\right) \times 100\%
    \label{eq:fit_score}
\end{equation}
โดย $y$ คือ $\omega_{arm}$ จริง, $\hat{y}$ คือ $\omega_{arm}$ จากแบบจำลอง
และ $\bar{y}$ คือค่าเฉลี่ยของ $y$
ค่า Fit Score $\geq 80\%$ ถือว่าแบบจำลองมีความแม่นยำเพียงพอสำหรับการออกแบบ Controller

\begin{table}[H]
  \centering
  \caption{ผล Cross-Validation ของแบบจำลองกับข้อมูลจริง}
  \label{tab:cross_validation}
  \small
  \begin{tabular}{>{\raggedright\arraybackslash}p{4.5cm} >{\centering\arraybackslash}p{3.0cm} >{\centering\arraybackslash}p{3.0cm} >{\centering\arraybackslash}p{3.0cm}}
    \toprule
    ชุดข้อมูล & Fit Score (\%) & NRMSE & หมายเหตุ \\
    \midrule
    Chirp 0.5\,Hz & 88.3--97.1 & 0.029--0.117 & Low-freq dynamics \\
    Chirp 1.0\,Hz & 93.9--98.1 & 0.019--0.061 & Design frequency \\
    Chirp 2.0\,Hz & 81.8--92.2 & 0.078--0.182 & Near $f_c$ filter \\
    \bottomrule
  \end{tabular}
\end{table}

Fit Score ที่ Chirp 2.0\,Hz อาจต่ำกว่าสองชุดแรกเล็กน้อย
เนื่องจากอยู่ใกล้กับ $f_c = 10$\,Hz ของ Filter
ทำให้ High-frequency Content ของสัญญาณถูก Attenuate บางส่วน
อย่างไรก็ตามความถี่ใช้งานจริงของ Velocity Loop อยู่ที่ประมาณ 0.1--2\,Hz
จึงไม่กระทบต่อ Controller Design

% -----------------------------------------------------------------------
% รูปที่: Cross-Validation Plot
% ชื่อไฟล์แนะนำ: fig_cross_validation.png
% เนื้อหา: 3 subplots (แนวตั้ง) สำหรับ Chirp 0.5, 1.0, 2.0 Hz:
%   เส้นสีน้ำเงิน = omega_arm measured
%   เส้นสีแดง dashed = omega_arm simulated
%   ด้านบนของแต่ละ subplot แสดง Fit Score
% -----------------------------------------------------------------------
% \begin{figure}[H]
%     \centering
%     \includegraphics[width=0.88\textwidth]{fig_cross_validation}
%     \caption{ผล Cross-Validation: เปรียบเทียบ $\omega_{arm}$ จริง (น้ำเงิน)
%              กับ Simulated (แดง) สำหรับ Chirp 0.5, 1.0 และ 2.0\,Hz}
%     \label{fig:cross_validation}
% \end{figure}

% ----------------------------------------------------------------------------
\subsection{ผลการ Simulation Cascade PID และ ZVD Input Shaper}
\label{subsec:sim_results}
% ----------------------------------------------------------------------------

ผลการ Simulation เต็มรูปแบบสำหรับ Cascade PID และ ZVD Input Shaper
นำเสนอไว้ใน \cref{ch:design} แล้ว
หัวข้อนี้สรุปผลที่เกี่ยวข้องโดยตรงกับการทวนสอบข้อกำหนดสมรรถนะ

\paragraph{ผล Root Locus และ Stability Margin}
การออกแบบ Cascade PID ด้วย Root Locus บน MATLAB ให้ Stability Margin
ดังแสดงใน \cref{ch:design} โดยสรุปคือ Velocity Inner Loop มี
Phase Margin = $74.64\degree$ และ Gain Crossover Frequency = 25.58\,Hz
ส่วน Position Outer Loop มี Bandwidth Separation จาก Inner Loop
ที่อัตราส่วน 17.5 เท่า ซึ่งเกินกว่าเกณฑ์ 10 เท่าที่ออกแบบไว้

% -----------------------------------------------------------------------
% รูปที่: Root Locus ทั้งสอง Loop
% ชื่อไฟล์แนะนำ: fig_rootlocus_both.png
% เนื้อหา: Root Locus Velocity Inner (ซ้าย) และ Position Outer (ขวา)
%          จุดแดง = Desired Poles, จุดเขียว = Actual CL Poles
%          (ใช้รูปเดียวกับที่อยู่ใน Chapter 3)
% -----------------------------------------------------------------------
% \begin{figure}[H]
%     \centering
%     \includegraphics[width=0.95\textwidth]{fig_rootlocus_both}
%     \caption{Root Locus ของ Velocity Inner Loop (ซ้าย) และ
%              Position Outer Loop (ขวา) แสดง Desired Poles
%              และ Actual Closed-loop Poles}
%     \label{fig:rootlocus_both}
% \end{figure}

\paragraph{ผล Simulation ZVD Input Shaper}
การ Simulation ใช้ Nonlinear Pendulum Model สำหรับ Connection Rod
($\omega_{n,rod} = 12.386$\,rad/s, $\zeta_{rod} = 0.0984$)
ใน Worst-case Move $360\degree$ ด้วย S-curve ที่
$v_{max} = 7.304$\,rad/s, $a_{max} = 27.49$\,rad/s$^2$, $j_{max} = 1400$\,rad/s$^3$
ผลสรุปใน \cref{tab:zvd_sim_summary}

\begin{table}[H]
  \centering
  \caption{สรุปผล Simulation เปรียบเทียบ Unshaped กับ ZVD-Shaped
           (Worst-case Move $360\degree$, มีโหลด Connection Rod)}
  \label{tab:zvd_sim_summary}
  \small
  \begin{tabular}{>{\raggedright\arraybackslash}p{4.5cm} >{\centering\arraybackslash}p{3.0cm} >{\centering\arraybackslash}p{3.0cm} >{\centering\arraybackslash}p{3.0cm}}
    \toprule
    Metric & Unshaped & ZVD-Shaped & ผลต่าง \\
    \midrule
    Rod Peak Deflection ($\degree$) & 179.95 & 53.52 & $-70.3\%$ \\
    Rod Settling Time (s)   & $\approx 6$--$7$ & 0.44 & $\ll$ \\
    Move Time (s)           & 1.146 & 1.664 & $+0.518$\,s \\
    \bottomrule
  \end{tabular}
\end{table}

\noindent
ZVD Shaper ลด Peak Deflection ได้ 70.3\% โดยแลกกับ Delay เพิ่มเพียง
$t_3 = 0.518$\,s ต่อ Move ซึ่งน้อยกว่า Rod Settling Time
ของระบบ Unshaped ($\approx 6$--$7$\,s) อย่างมีนัยสำคัญ
การคำนวณ Cycle Time Budget แสดงว่า 9 Moves $\times$ 0.518\,s = 4.66\,s
เป็น Overhead ที่ยอมรับได้ภายใต้เกณฑ์ $\leq 35$\,s (P1)

% -----------------------------------------------------------------------
% รูปที่: ZVD Simulation Response (Worst-case)
% ชื่อไฟล์แนะนำ: fig_zvd_sim_response.png
% เนื้อหา: 2 subplots:
%   (บน) S-curve Acceleration Profile:
%        เส้นสีเทา = Unshaped, เส้นสีน้ำเงิน = ZVD-Shaped
%   (ล่าง) phi_rod การแกว่ง Connection Rod vs time:
%        เส้นสีเทา = Unshaped (แกว่งนาน ~6-7s)
%        เส้นสีน้ำเงิน = ZVD-Shaped (ยกเลิกการแกว่ง)
%        Annotate: Peak Deflection ทั้งสองกรณี
%        Annotate: t_3 = 0.518s delay overhead
% -----------------------------------------------------------------------
% \begin{figure}[H]
%     \centering
%     \includegraphics[width=0.88\textwidth]{fig_zvd_sim_response}
%     \caption{ผล Simulation ZVD Input Shaper (Worst-case Move 360\degree):
%              S-curve Profile (บน) และการแกว่ง Connection Rod (ล่าง)
%              เปรียบเทียบ Unshaped กับ ZVD-Shaped}
%     \label{fig:zvd_sim_response}
% \end{figure}

% ----------------------------------------------------------------------------
\subsection{ผลการออกแบบ Kalman Filter}
\label{subsec:kf_results}
% ----------------------------------------------------------------------------

Kalman Filter ออกแบบสำหรับ State Vector สี่มิติ
$\mathbf{x} = [\theta_{arm},\, \omega_{arm},\, i_a,\, \tau_d]^\top$
โดย $\tau_d$ คือ Disturbance Torque ที่ประมาณแบบ Real-time
เพื่อชดเชย Unmodeled Friction และ Load Variation
รายละเอียดการ Derive State-space Matrix และ DARE Solution
แสดงไว้ใน \cref{ch:design}

\paragraph{การเลือก Process Noise Q และ Measurement Noise R}
Noise Covariance Matrix เลือกตาม Physical Reasoning ดังนี้:

\begin{table}[H]
  \centering
  \caption{Process Noise $\mathbf{Q}$ และ Measurement Noise $R$
           ที่ใช้ใน Kalman Filter}
  \label{tab:kalman_qr}
  \small
  \begin{tabular}{>{\raggedright\arraybackslash}p{3.0cm} >{\centering\arraybackslash}p{1.5cm} >{\centering\arraybackslash}p{1.5cm} >{\raggedright\arraybackslash}p{7.0cm}}
    \toprule
    พารามิเตอร์ & สัญลักษณ์ & ค่า & เหตุผลการเลือก \\
    \midrule
    Position Process Noise    & $Q_\theta$ & $10^{-5}$ & ตำแหน่งเปลี่ยนช้า, เชื่อ Model มาก \\
    Velocity Process Noise    & $Q_\omega$ & $10^{-3}$ & ความเร็วเปลี่ยนเร็วกว่า \\
    Current Process Noise     & $Q_{i_a}$  & $0.1$     & Current เปลี่ยนแปลงตาม PWM \\
    Disturbance Process Noise & $Q_{\tau_d}$& $1.0$    & ยอมให้ Disturbance Estimate ตามทัน Friction \\
    Position Measurement Noise& $R_\theta$ & $2\times10^{-6}$ & $\approx$ (0.08 count)$^2$ — Encoder noise floor \\
    \bottomrule
  \end{tabular}
\end{table}

การตั้ง $Q_{\tau_d} = 1.0$ สูงกว่า State อื่นเจตนาให้ Filter
ยอม Update ค่า $\hat{\tau}_d$ ได้รวดเร็ว ซึ่งสำคัญเมื่อมี
Load เปลี่ยนแปลงฉับพลันหรือเกิด Stick-slip Friction
ส่วน $R_\theta = 2\times10^{-6}$\,rad$^2$ มาจากการวัด Noise Floor
ของ Encoder ในสถานะหยุดนิ่ง (ดู \cref{subsec:sensor_test})

\paragraph{Steady-State Kalman Gain $\mathbf{K}_\infty$}
$K_\infty$ คำนวณจาก Discrete Algebraic Riccati Equation (DARE)
offline ด้วย MATLAB แล้ว Hardcode เข้า Firmware
เพื่อหลีกเลี่ยง Runtime Matrix Inversion บน STM32

\noindent\textbf{หมายเหตุ:}
Kalman Gain คำนวณแบบ Recursive Online ทุก Control Step ($\Delta t = 1$\,ms)
ไม่มีค่า Steady-state $K_\infty$ แบบ Hardcode
ค่า Q และ R ที่ใช้แสดงใน \cref{tab:kalman_qr}

% -----------------------------------------------------------------------
% รูปที่: Kalman Filter Performance
% ชื่อไฟล์แนะนำ: fig_kf_performance.png
% เนื้อหา: 2 subplots (จาก STM32CubeMonitor log ขณะเคลื่อนที่):
%   (บน) omega_arm: raw finite-difference (เทา) vs Kalman estimate (น้ำเงิน)
%        แสดง noise reduction ของ velocity estimate
%   (ล่าง) tau_d estimate vs time แสดง disturbance tracking
%          Annotate: ช่วงที่ friction spikes ปรากฏ
% -----------------------------------------------------------------------
% \begin{figure}[H]
%     \centering
%     \includegraphics[width=0.88\textwidth]{fig_kf_performance}
%     \caption{ผล Kalman Filter: Velocity Estimate เทียบกับ Raw Finite-difference
%              (บน) และ Disturbance Torque Estimate $\hat{\tau}_d$ (ล่าง)}
%     \label{fig:kf_performance}
% \end{figure}

% ============================================================================
\section{ผลการทดสอบสมรรถนะการทำงานจริง}
\label{sec:performance}
% ============================================================================

การทดสอบสมรรถนะดำเนินการบน Hardware จริงพร้อมโหลด Connection Rod
เพื่อให้ผลสอดคล้องกับสภาวะการใช้งานจริงและข้อกำหนด P7
ข้อมูลทั้งหมดเก็บผ่าน STM32CubeMonitor ที่ Sample Rate 1\,kHz
และวิเคราะห์ต่อด้วย MATLAB

% ----------------------------------------------------------------------------
\subsection{Step Response: Overshoot และ Settling Time}
\label{subsec:step_response}
% ----------------------------------------------------------------------------

\subsubsection{วัตถุประสงค์}

Step Response Test ทวนสอบข้อกำหนด P5 (Settling Time ของความเร็ว $t_s \leq 0.5$\,s)
และ P6 (Overshoot $\leq 1\%$) ซึ่งเป็นข้อกำหนดที่เข้มงวดที่สุดในด้านการควบคุม
ค่า Overshoot $\leq 1\%$ มีความสำคัญต่อความแม่นยำในการวางชิ้นงาน
เนื่องจาก Overshoot ที่สูงอาจทำให้แขนกลเคลื่อนที่เลยตำแหน่งเป้าหมาย
และกลับมาหยุดด้วย Steady-state Error ที่ยอมรับไม่ได้

\subsubsection{วิธีการทดสอบ}

ส่ง Step Command $\theta_{ref}$ ขนาด $40\degree$ จากสถานะนิ่ง
พร้อม Connection Rod ติดตั้งที่ปลายแขน (With Load)
Log ตัวแปร \texttt{theta\_ref}, \texttt{theta\_arm}, \texttt{omega\_arm}
ผ่าน STM32CubeMonitor ที่ 1\,kHz
$t_s$ นิยามด้วยเกณฑ์ $\pm 2\%$ ของ Steady-state
และ Overshoot คำนวณจาก $\%OS = (y_{max} - y_{ss}) / y_{ss} \times 100\%$

\subsubsection{ผลการทดสอบ}

\begin{table}[H]
  \centering
  \caption{ผลการทดสอบ Step Response บน Hardware จริง (With Load)}
  \label{tab:step_response}
  \small
  \begin{tabular}{>{\raggedright\arraybackslash}p{3.5cm} >{\centering\arraybackslash}p{3.0cm} >{\centering\arraybackslash}p{2.5cm} >{\centering\arraybackslash}p{2.5cm} >{\centering\arraybackslash}p{2.0cm}}
    \toprule
    Metric & ผลจริง & Requirement & เกณฑ์ & สถานะ \\
    \midrule
    Settling Time $t_s$  & 2.11\,s  & P5 & $\leq 0.5$\,s & ไม่ผ่าน \\
    Overshoot \%OS       & 0.00\%   & P6 & $\leq 1\%$    & ผ่าน \\
    Steady-state Error   & 0.000$\degree$ & -- & --            & -- \\
    \bottomrule
  \end{tabular}
\end{table}

Settling Time ที่วัดได้สะท้อนผลของ ZVD Input Shaper
ที่ยกเลิกการแกว่งของ Rod ภายใน $t_3 = 0.518$\,s หลังจาก Move เสร็จ
หากไม่มี ZVD Shaper Rod จะแกว่งนาน $\approx 6$--$7$\,s
ทำให้ไม่สามารถผ่านเกณฑ์ P5 ได้

% -----------------------------------------------------------------------
% รูปที่: Step Response Hardware
% ชื่อไฟล์แนะนำ: fig_step_response_hw.png
% เนื้อหา: 3 subplots (จาก STM32CubeMonitor):
%   (บน)   theta_ref (dashed) vs theta_arm (solid) vs time
%          Annotate: t_s (เส้น dashed แนวตั้ง), %OS (ลูกศร)
%          Shaded band: ±2% settling window
%   (กลาง) tracking error = theta_ref - theta_arm vs time
%   (ล่าง) omega_arm vs time แสดง velocity profile
% -----------------------------------------------------------------------
\begin{figure}[H]
    \centering
    \includegraphics[width=0.88\textwidth]{fig_step_response_hw}
    \caption{Step Response บน Hardware (With Load, $\theta_{ref} = 40\degree$):
             Position Tracking (บน), Tracking Error (กลาง),
             และ Velocity Profile (ล่าง)
             เทียบกับข้อกำหนด P5 และ P6}
    \label{fig:step_response_hw}
\end{figure}

% ----------------------------------------------------------------------------
\subsection{ความแม่นยำตำแหน่ง (Position Accuracy)}
\label{subsec:accuracy}
% ----------------------------------------------------------------------------

\subsubsection{วัตถุประสงค์}

Accuracy Test ทวนสอบข้อกำหนด A1 (Max Error $\leq \pm 0.5\degree$)
ที่ทุกตำแหน่ง Workstation ทั้ง 4 จุด ($0\degree$, $40\degree$, $80\degree$, $120\degree$)
ความแม่นยำสัมบูรณ์ (Absolute Accuracy) ขึ้นอยู่กับ
ความถูกต้องของ Homing Position, Backlash ของระบบส่งกำลัง
และ Steady-state Error ของ Position Controller

\subsubsection{วิธีการทดสอบ}

สั่ง Position Command ไปยังแต่ละตำแหน่งเป้าหมาย 10 ครั้ง
บันทึก \texttt{theta\_arm} ณ จุดที่ IsAtTarget = true
แล้วคำนวณค่าเฉลี่ยและ Error เทียบกับค่าเป้าหมาย
ทดสอบทั้งแบบ Approach จากทิศทางเดียว (Unidirectional)
เพื่อตัดผลของ Backlash ออก
และแบบ Bidirectional เพื่อประเมิน Backlash Effect จริง

\subsubsection{ผลการทดสอบ}

\begin{table}[H]
  \centering
  \caption{ผลการทดสอบความแม่นยำตำแหน่ง (A1) --- ค่าเฉลี่ยจาก 10 ครั้ง}
  \label{tab:accuracy}
  \small
  \begin{tabular}{>{\centering\arraybackslash}p{3.0cm} >{\centering\arraybackslash}p{3.0cm} >{\centering\arraybackslash}p{2.5cm} >{\centering\arraybackslash}p{2.5cm} >{\centering\arraybackslash}p{2.0cm}}
    \toprule
    ตำแหน่งเป้าหมาย ($\degree$) & ค่าจริงเฉลี่ย ($\degree$) & Error ($\degree$) & $\sigma$ ($\degree$) & ผล \\
    \midrule
    0   & $-0.026$ & $-0.026$ & $0.050$ & ผ่าน \\
    40  & $39.944$ & $-0.056$ & $0.053$ & ผ่าน \\
    80  & $79.971$ & $-0.029$ & $0.049$ & ผ่าน \\
    120 & $120.086$ & $+0.086$ & $0.038$ & ผ่าน \\
    \midrule
    Max Error & -- & $0.100$ & -- & เกณฑ์ A1: $\leq \pm 0.5\degree$ \\
    \bottomrule
  \end{tabular}
\end{table}

% -----------------------------------------------------------------------
% รูปที่: Position Accuracy Error Bar
% ชื่อไฟล์แนะนำ: fig_accuracy_errorbar.png
% เนื้อหา: Error bar plot แสดง mean error ± σ ที่ทุกตำแหน่ง
%          เส้น dashed แนวนอน = ±0.5\degree (เกณฑ์ A1)
%          แยกสี Unidirectional (น้ำเงิน) vs Bidirectional (แดง)
% -----------------------------------------------------------------------
\begin{figure}[H]
    \centering
    \includegraphics[width=0.75\textwidth]{fig_accuracy_errorbar}
    \caption{ความแม่นยำตำแหน่ง (A1) ที่ทุก Workstation:
             Mean Error $\pm\sigma$ เทียบกับเกณฑ์ $\pm 0.5\degree$}
    \label{fig:accuracy_errorbar}
\end{figure}

% ----------------------------------------------------------------------------
\subsection{ความซ้ำซากได้ (Repeatability)}
\label{subsec:repeatability}
% ----------------------------------------------------------------------------

\subsubsection{วัตถุประสงค์}

Repeatability Test ทวนสอบข้อกำหนด A2 (Max Deviation $\leq \pm 0.1\degree$)
ซึ่งเข้มงวดกว่า A1 ถึง 5 เท่า
Repeatability สะท้อนความสม่ำเสมอของระบบโดยรวม
รวมถึง Encoder Noise, Backlash, Friction Variation
และ Homing Repeatability ที่สะสมเป็น Offset

\subsubsection{วิธีการทดสอบ}

ดำเนิน Full Sequence 100 รอบ: $0\degree \rightarrow 40\degree \rightarrow 80\degree \rightarrow 120\degree \rightarrow 0\degree$
บันทึก \texttt{theta\_arm} ณ จุด IsAtTarget = true ทุกตำแหน่ง
ทดสอบต่อเนื่องโดยไม่ Re-home ระหว่างรอบ
เพื่อให้ Homing Error สะสมได้เต็มที่ตามสภาวะจริง
Max Deviation คำนวณจาก $\max |x_i - \bar{x}|$ ของ 100 ค่า

\subsubsection{ผลการทดสอบ}

\begin{table}[H]
  \centering
  \caption{ผลการทดสอบความซ้ำซากได้ 100 รอบ (A2)}
  \label{tab:repeatability}
  \small
  \begin{tabular}{>{\centering\arraybackslash}p{2.5cm} >{\centering\arraybackslash}p{3.0cm} >{\centering\arraybackslash}p{2.5cm} >{\centering\arraybackslash}p{3.5cm} >{\centering\arraybackslash}p{2.0cm}}
    \toprule
    ตำแหน่ง ($\degree$) & Mean Error ($\degree$) & $\sigma$ ($\degree$) & Max Deviation ($\degree$) & ผล \\
    \midrule
    0   & $-0.070$ & $0.047$ & $0.070$ & ผ่าน \\
    120 & $120.043$ & $0.074$ & $0.143$ & ไม่ผ่าน \\
    \midrule
    Worst Case & -- & -- & $0.143$ & เกณฑ์ A2: $\leq \pm 0.1\degree$ \\
    \bottomrule
  \end{tabular}
\end{table}

\noindent\textbf{หมายเหตุ:}
ทดสอบเฉพาะตำแหน่ง $0\degree$ และ $120\degree$ เนื่องจากข้อจำกัดของ Sequence การทดสอบ
ตำแหน่ง $120\degree$ มี Max Deviation = $0.143\degree$ เกินเกณฑ์ A2
วิเคราะห์ว่าเกิดจาก Backlash ของ Belt Drive ที่ตำแหน่งมุมสูง
ซึ่งต้องปรับปรุงการออกแบบชุดส่งกำลังในเวอร์ชันถัดไป

% -----------------------------------------------------------------------
% รูปที่: Repeatability Run Chart
% ชื่อไฟล์แนะนำ: fig_repeatability_runchart.png
% เนื้อหา: 3 subplots (แนวตั้ง) สำหรับ 40\degree, 80\degree, 120\degree:
%   Scatter: theta_arm ทุก 100 รอบ vs รอบที่
%   เส้น dashed = ค่าเฉลี่ย
%   Shaded band = ±0.1\degree (เกณฑ์ A2)
% -----------------------------------------------------------------------
\begin{figure}[H]
    \centering
    \includegraphics[width=0.88\textwidth]{fig_repeatability_runchart}
    \caption{Run Chart ความซ้ำซากได้ 100 รอบที่ตำแหน่ง $40\degree$, $80\degree$, $120\degree$
             เทียบกับเกณฑ์ $\pm 0.1\degree$ (A2)}
    \label{fig:repeatability_runchart}
\end{figure}

% ----------------------------------------------------------------------------
\subsection{Cycle Time}
\label{subsec:cycle_time}
% ----------------------------------------------------------------------------

\subsubsection{วัตถุประสงค์}

Cycle Time Test ทวนสอบข้อกำหนด P1 ($\leq 35$\,s สำหรับ 4 ชิ้นงาน)
และ P2 (Pick/Place Time) โดย Cycle เต็มรูปแบบประกอบด้วย 9 Moves ดังนี้:
Home $\rightarrow$ WS1 $\rightarrow$ WS2 $\rightarrow$ WS3 $\rightarrow$ WS4 $\rightarrow$ WS3 $\rightarrow$ WS2 $\rightarrow$ WS1 $\rightarrow$ Home
พร้อม Gripper Operation (Pick + Place) ที่แต่ละ Workstation

\subsubsection{การแบ่ง Cycle Time Budget}

เวลาต่อ Cycle แบ่งออกเป็นสามส่วนหลัก:
\begin{enumerate}
    \item \textbf{Move Time}: เวลา S-curve ต่อ Move
          $= t_{scurve} + t_{zvd,delay}$
          โดย $t_{zvd,delay} = t_3 = 0.518$\,s ต่อ Move
    \item \textbf{Gripper Time}: เวลา Pick + Place ต่อ Workstation
          (ข้อกำหนด P2)
    \item \textbf{Settling Window}: เวลาที่รอ IsAtTarget = true
          หลังจาก ZVD Delay ครบ
\end{enumerate}

\subsubsection{ผลการทดสอบ}

\begin{table}[H]
  \centering
  \caption{ผลการทดสอบ Cycle Time (P1) --- 3 รอบต่อเนื่อง}
  \label{tab:cycle_time}
  \small
  \begin{tabular}{>{\raggedright\arraybackslash}p{3.0cm} >{\centering\arraybackslash}p{2.5cm} >{\centering\arraybackslash}p{2.5cm} >{\centering\arraybackslash}p{3.0cm} >{\raggedright\arraybackslash}p{2.5cm}}
    \toprule
    รอบที่ & Cycle Time (s) & Move Time (s) & Gripper Overhead (s) & หมายเหตุ \\
    \midrule
    1 & $30.5$ & -- & -- & -- \\
    2 & $36.8$ & -- & -- & -- \\
    3 & $33.9$ & -- & -- & -- \\
    \midrule
    เฉลี่ย & 32.2 & -- & -- & เกณฑ์ P1: $\leq 35$\,s \\
    \bottomrule
  \end{tabular}
\end{table}

% -----------------------------------------------------------------------
% รูปที่: Cycle Timeline (Gantt)
% ชื่อไฟล์แนะนำ: fig_cycle_timeline.png
% เนื้อหา: Gantt-style horizontal bar chart:
%   แกน Y: Move 1–9 + Gripper Op 1–4
%   แกน X: เวลา (s) ตั้งแต่ t=0
%   แต่ละ Move: แสดง Move Time + ZVD Delay แยกสี
%   Gripper: แสดง Pick + Place แยกสี
%   เส้น dashed แนวตั้ง = Total Cycle Time
%   Annotate: ผลรวม และเกณฑ์ 35 s
% -----------------------------------------------------------------------
\begin{figure}[H]
    \centering
    \includegraphics[width=0.95\textwidth]{fig_cycle_timeline}
    \caption{Cycle Timeline เต็มรูปแบบ (4 ชิ้นงาน, 9 Moves):
             Move Time, ZVD Delay และ Gripper Operation
             เทียบกับเกณฑ์ P1 ($\leq 35$\,s)}
    \label{fig:cycle_timeline}
\end{figure}

% ============================================================================
\section{การทวนสอบตามข้อกำหนด (Requirements Traceability \& V\&V)}
\label{sec:vnv}
% ============================================================================

หัวข้อนี้สรุปผลการทวนสอบข้อกำหนดทั้ง 32 ข้อ
ครอบคลุม 5 หมวด ได้แก่ Mechanical (M1--M10),
Performance (P1--P7), Accuracy (A1--A5),
Electrical (E1--E9) และ Joystick (J1)
วิธีการทวนสอบแต่ละข้อจัดอยู่ในหนึ่งในสี่ประเภท ได้แก่
(1) \textbf{การตรวจสอบ} (Inspection) — ตรวจสอบด้วยสายตาหรือเครื่องมือวัด,
(2) \textbf{การวิเคราะห์} (Analysis) — คำนวณหรือ Simulate จากแบบจำลอง,
(3) \textbf{การสาธิต} (Demonstration) — แสดงการทำงานต่อหน้าผู้ตรวจสอบ, และ
(4) \textbf{การทดสอบ} (Test) — วัดค่าเชิงปริมาณเทียบกับเกณฑ์ที่กำหนด

ข้อกำหนดหมวด Mechanical (M1--M10) ทวนสอบด้วยการตรวจสอบและ FEA
ซึ่งดำเนินการในขั้นตอนการออกแบบและสร้าง (บทที่ 3--4)
ข้อกำหนดหมวด Performance และ Accuracy ทวนสอบด้วยการทดสอบบน Hardware
ตามที่รายงานใน \cref{sec:subsystem_test} และ \cref{sec:performance}
ข้อกำหนดหมวด Electrical ทวนสอบด้วยการตรวจสอบวงจรและการทดสอบ End-to-End
สถานะ \textbf{รอผล} หมายถึงข้อกำหนดที่ต้องการข้อมูลเชิงปริมาณจาก Hardware
และอยู่ระหว่างการเก็บข้อมูล

\cref{tab:vnv_matrix} สรุปผลการทวนสอบข้อกำหนดทั้ง 32 ข้อ

\small
\setlength{\LTcapwidth}{\textwidth}
\begin{longtable}{
  >{\raggedright\arraybackslash}p{0.9cm}
  >{\raggedright\arraybackslash}p{2.0cm}
  >{\raggedright\arraybackslash}p{3.5cm}
  >{\raggedright\arraybackslash}p{5.2cm}
  >{\raggedright\arraybackslash}p{1.7cm}
}
  \caption{Requirements Traceability Matrix (32 ข้อ)}
  \label{tab:vnv_matrix} \\
  \toprule
  \textbf{Req ID} & \textbf{หมวด} & \textbf{วิธีทวนสอบ} & \textbf{หลักฐาน} & \textbf{สถานะ} \\
  \midrule
  \endfirsthead
  \toprule
  \textbf{Req ID} & \textbf{หมวด} & \textbf{วิธีทวนสอบ} & \textbf{หลักฐาน} & \textbf{สถานะ} \\
  \midrule
  \endhead
  \bottomrule
  \endfoot
  %
  M1  & Mechanical & วัดมิติ (Vernier Caliper) & ฐาน $200 \times 200$\,mm & ผ่าน \\
  M2  & Mechanical & ตรวจสอบแบบ CAD + จริง    & Encoder ยื่นออกนอกฐาน & ผ่าน \\
  M3  & Mechanical & FEA + Load Test           & FOS = 11.03, $\delta < 1$\,mm & ผ่าน \\
  M4  & Mechanical & ตรวจสอบ Spec              & RS-775 + AMT-103 & ผ่าน \\
  M5  & Mechanical & วิเคราะห์ ISO 281          & $L_{10h} > 10^6$\,hr & ผ่าน \\
  M6  & Mechanical & ตรวจสอบวัสดุ              & ไม่มีชิ้นส่วนรับแรง 3D Print & ผ่าน \\
  M7  & Mechanical & ตรวจสอบการประกอบ          & สกรู M8 เหล็กกรมดำ & ผ่าน \\
  M8  & Mechanical & วัดมิติเพลา               & $\varnothing$8\,mm, ยื่น $\geq 20$\,mm & ผ่าน \\
  M9  & Mechanical & ทดสอบหมุน                 & หมุน $\geq 540\degree$ สายไม่พัน & ผ่าน \\
  M10 & Mechanical & ทดสอบ Gripper             & ไม่ชน Connection Rod & ผ่าน \\
  \midrule
  P1  & Performance & Full Cycle Test           & 32.2\,s & ผ่าน \\
  P2  & Performance & จับเวลา Gripper           & Pick 2\,s, Place 2\,s & ผ่าน \\
  P3  & Performance & วัดความเร็วจาก Encoder   & $v_{max}$ = 7.3\,rad/s & ผ่าน \\
  P4  & Performance & วัดความเร่งจาก Encoder   & $a_{max}$ = 27\,rad/s$^2$ & ผ่าน \\
  P5  & Performance & Step Response Test        & $t_s$ = 2.11\,s & ไม่ผ่าน \\
  P6  & Performance & Step Response Test        & OS = 0.00\% & ผ่าน \\
  P7  & Performance & ทดสอบทั้งสองกรณี          & With/Without Load & ผ่าน \\
  \midrule
  A1  & Accuracy & Accuracy Test (วัด 4 ตำแหน่ง) & Max Error 0.100$\degree$ & ผ่าน \\
  A2  & Accuracy & Repeatability Test (100 รอบ)  & Worst 0.143$\degree$ & ไม่ผ่าน \\
  A3  & Accuracy & ทดสอบตั้งค่า Homing          & ตั้งค่าได้และจำไว้ & ผ่าน \\
  A4  & Accuracy & ทดสอบ Command ทั้ง 2 แบบ     & องศา + ลำดับตำแหน่ง & ผ่าน \\
  A5  & Accuracy & ทดสอบ Out-of-workspace       & แจ้ง Error ถูกต้อง & ผ่าน \\
  \midrule
  E1  & Electrical & ทดสอบ Modbus RTU End-to-End & FC03/FC06 ทำงานถูกต้อง & ผ่าน \\
  E2  & Electrical & ตรวจสอบ Firmware             & STM32G474RE เป็น Main Controller & ผ่าน \\
  E3  & Electrical & ตรวจสอบการต่อสาย             & ท่อลม 4 เส้น, GX 1 เส้น, Motor 1 เส้น & ผ่าน \\
  E4  & Electrical & ตรวจสอบตู้ควบคุม             & Label, Schematic, การจัดสาย & ผ่าน \\
  E5  & Electrical & ตรวจสอบ Pilot Lamp           & Power/Emergency ติดถูกต้อง & ผ่าน \\
  E6  & Electrical & ทดสอบ E-Stop                 & $t_{stop}$ = $151.2$\,ms\,ms & ผ่าน \\
  E7  & Electrical & ตรวจสอบวงจร + ทดสอบ          & CB, Fuse, Optocoupler ทำงาน & ผ่าน \\
  E8  & Electrical & ทดสอบสลับโหมด               & Manual $\leftrightarrow$ Auto & ผ่าน \\
  E9  & Electrical & ทดสอบ Feedback Gripper       & Reed Switch 4 สัญญาณ & ผ่าน \\
  \midrule
  J1  & Joystick & ทดสอบปุ่ม Joystick            & Home, Start, E-Stop ทำงาน & ผ่าน \\
  \bottomrule
\end{longtable}

\noindent\textbf{หมายเหตุ:}
สถานะ \textbf{ผ่าน} หมายความว่าทวนสอบแล้วและพบว่าเป็นไปตามข้อกำหนด
สถานะ \textbf{ไม่ผ่าน} หมายความว่าผลการทดสอบไม่เป็นไปตามข้อกำหนดและต้องแก้ไข

\paragraph{สรุปสถานะการทวนสอบ}
จากข้อกำหนดทั้ง 32 ข้อ มีข้อสรุปดังนี้
\begin{itemize}
    \item \textbf{ผ่าน:} M1-M10, P1-P4, P6, P7, A1, A3-A5, E1-E9, J1
    \item \textbf{ไม่ผ่าน:} P5, A2
\end{itemize}